% ID: SEL-D0034-OB-E01
% ESTADO: validada. S035, cotejada con examen y criterios.
\textbf{a)} El recinto queda por encima de $y=2x+2$ y
$y=5x-4$, por debajo de $y=2x+5$ y a la derecha de $x=0$.
En el eje vertical van de $(0,2)$ a $(0,5)$. La recta
$y=5x-4$ corta a las dos paralelas en $(2,6)$ y $(3,11)$,
respectivamente.
\begin{center}
\setlength{\unitlength}{5mm}
\begin{picture}(9,14)(-3,-1)
\qbezier(0,0)(2.5,0)(5,0)\put(5.2,0){$x$}
\qbezier(0,0)(0,6)(0,12)\put(0,12.2){$y$}
\qbezier(0,2)(1,4)(2,6)
\qbezier(2,6)(2.5,8.5)(3,11)
\qbezier(3,11)(1.5,8)(0,5)
\put(-0.2,2){\makebox(0,0)[r]{$(0,2)$}}
\put(-0.2,5){\makebox(0,0)[r]{$(0,5)$}}
\put(2.2,6){$(2,6)$}
\put(3.2,11){$(3,11)$}
\end{picture}
\end{center}
\textbf{b)} Escribimos $F=-3(y-2x)$. Como $2\le y-2x\le5$,
el máximo es $\boxed{-6}$ en todo el lado $y=2x+2$,
$0\le x\le2$. El mínimo es $\boxed{-15}$ en todo el lado
$y=2x+5$, $0\le x\le3$. En cada caso la igualdad fija la
recta, y las restricciones restantes fijan el segmento.
